\section{The algebraic locus and the propagation criteria}\label{sec:propagation}

\subsection{The algebraic locus is closed, and what that costs}
\label{ssec:properness}

Everything so far has been pointwise: a class at a point, a criterion at a
point, a cycle at a point, and in \Cref{cor:exactform} an implication to be
proved at every point. The propagation question of \Cref{q:residual} is about
the set
\[
  \Sigma \;=\; \bigl\{\,s\in\cD_{n,n} \;:\;
    \omega_{s}\ \text{is algebraic on}\ A_{s}\,\bigr\},
\]
and nothing has yet been said about what kind of set that is. It turns out to
be a very restricted kind, and the restriction changes the shape of the
remaining problem: instead of a class to construct at every point, one is left
with a bound to prove along one orbit, or a rank to compute at one point.
The mechanism is that the Hilbert scheme is proper, so each piece of $\Sigma$
cut out by bounded data is Zariski closed.\footnote{The mechanism is not new.
That the locus where a fixed Hodge class becomes algebraic is a countable
union of closed subvarieties is the standard companion to the statement that
the Hodge locus itself is a countable union of closed subvarieties, for which
see \cite[Chapter~5]{Voi03} and, in the form we use in \Cref{thm:cdk},
\cite{CDK95}. What is specific to the present situation is that the class
whose algebraic locus is being described has been pinned down completely by
\Cref{thm:hdgcount} and written out by \Cref{thm:explicitgens}, that the locus
is known to be nonempty in every family by \Cref{thm:everyfamily}, and that
the group acting has been identified. Those three inputs are what turn a soft
statement into a criterion one can try to verify.}

\begin{setupenv}[The quotient, and the universal family]\label{setup:shimura}
Keep \Cref{setup:family}. Fix an integer $N\ge3$ and let $\Gamma$ be the
subgroup of $\SU(V,H)(\QQ)$ preserving the lattice $L\subset V$ used to build
the family and acting trivially on $L/NL$. Write
\[
  p\colon\cD_{n,n}\longrightarrow\cS=\Gamma\backslash\cD_{n,n} .
\]
By Baily and Borel \cite{BB66} the quotient $\cS$ is a smooth quasi-projective
complex variety, and it is irreducible because $\cD_{n,n}$ is connected. Since
$N\ge3$ the level structure rigidifies, so $\cS$ carries a universal family
\[
  f\colon\cA\longrightarrow\cS,
\]
and the relative polarisation $3\eta$ is relatively very ample, which fixes a
relative projective embedding and hence a notion of degree in the fibres.
Write $G=\SU(V,H)$, an algebraic group over $\QQ$, so that
$G(\RR)$ acts transitively on $\cD_{n,n}$ and $\Gamma\subset G(\QQ)$.
\end{setupenv}

Two small facts have to be recorded before the main one, because both are used
without comment later and neither is quite automatic.

\begin{lemma}[The section descends]\label{lem:omegadescends}
The section $\omega$ of \Cref{setup:family} is $\Gamma$-invariant, hence
descends to a flat section of $R^{2n}f_{*}\QQ$ over $\cS$ that is of type
$(n,n)$ at every point.
\end{lemma}

\begin{proof}
The Weil line is $\HW=\bigwedge^{2n}_{K}V$, and an element of $\SU(V,H)$ is by
definition a $K$-linear automorphism of $V$ of $K$-determinant $1$. Its action
on the top $K$-exterior power is multiplication by that determinant, so
$\Gamma$ acts trivially on $\HW$ and in particular fixes $\omega$. Flatness is
then immediate, and the Hodge type at every point is
\Cref{prop:weilline}(ii).
\end{proof}

\begin{lemma}[One component is enough]\label{lem:onecomponent}
Let $(A,\iota,\eta)$ be of $(K,-1,n)$-Weil type with
$\iota(\cO_{K})\subseteq\End(A)$, and let $\HW\subset H^{2n}(A,\QQ)$ be its
Weil line. If one nonzero element of $\HW$ is algebraic, every element of
$\HW$ is algebraic.\footnote{This is the step compressed into the last
sentence of the proof of \Cref{cor:whatisleft}, written out because the
witness $\tau$ has to be chosen and because the choice is the same one the
Lean theorem \texttt{irrational\_witness\_exists} certifies. The determinant
that makes the two classes independent is $b(w_{1}^{2}+dw_{2}^{2})$ for
$\tau^{2n}=a+b\sqrt{-d}$ and $w=w_{1}+w_{2}\sqrt{-d}$, and that is the content
of \texttt{weil\_line\_spanned}.}
\end{lemma}

\begin{proof}
Among the elements $\tau_{m}=m+\sqrt{-d}$ of $\cO_{K}$ with $1\le m\le7$ the
ratios $\tau_{m}/\bbar{\tau}_{m}$ are pairwise distinct, since
$\tau_{m}\bbar{\tau}_{m'}$ is rational only when $m=m'$. The group of roots of
unity in an imaginary quadratic field has order $2$, $4$ or $6$, so some $m$
in that range has $\tau=\tau_{m}$ with $\tau/\bbar\tau$ not a root of unity,
and then $\tau^{2n}\ne\bbar{\tau}^{2n}$, that is $\tau^{2n}\notin\QQ$, for
every $n\ge1$.\footnote{The exceptional witnesses are real and are the reason
the argument names a range rather than taking $m=1$: for $d=1$ the ratio
attached to $1+\sqrt{-1}$ is a primitive fourth root of unity and for $d=3$
the ratio attached to $1+\sqrt{-3}$ is a primitive sixth root of unity, so at
those two values of $d$ the naive witness fails for the divisible $n$. The
Lean theorems \texttt{bad\_witness} and \texttt{bad\_witness\_exceptions} of
\Cref{sec:lean} record exactly which pairs fail and that no others do in the
range searched.}

Since $\tau\in\cO_{K}$, the element $\iota(\tau)$ is an endomorphism of $A$,
so $\iota(\tau)^{*}$ on $H^{2n}(A,\QQ)$ is induced by the graph
correspondence and therefore carries algebraic classes to algebraic classes.
On $H^{1}$ it is multiplication by $\tau$, so on
$\bigwedge^{2n}_{K}V$ it is multiplication by $\tau^{2n}$. As
$\tau^{2n}\notin\QQ$ and $\HW$ is a one dimensional $K$-vector space, the pair
$w,\;\tau^{2n}w$ is a $\QQ$-basis of $\HW$ for every $w\ne0$. If $w$ is
algebraic then so is $\iota(\tau)^{*}w=\tau^{2n}w$, and the two together span.
\end{proof}

\begin{theorem}[The algebraic locus is a countable union of closed
subvarieties]\label{thm:closedlocus}
Keep \Cref{setup:shimura}. For a finite tuple $\underline{P}=(P_{1},\dots
,P_{m})$ of Hilbert polynomials and a tuple $\underline{q}\in\QQ^{m}$ put
\begin{align*}
  \cH_{\underline{P}}
  &\;=\;
  \Hilb^{P_{1}}(\cA/\cS)\times_{\cS}\cdots\times_{\cS}\Hilb^{P_{m}}(\cA/\cS),
  \\[2pt]
  \Sigma_{\underline{P},\underline{q}}
  &\;=\;
  \mathrm{pr}\Bigl(\bigl\{\,\underline{Z}\in\cH_{\underline{P}}
    \;:\;\textstyle\sum_{i}q_{i}[Z_{i}]=\omega\,\bigr\}\Bigr)
  \;\subseteq\;\cS .
\end{align*}
Then each $\Sigma_{\underline{P},\underline{q}}$ is Zariski closed in
$\cS$,\footnote{The bound on the Hilbert polynomials is what the statement
lives on, and it cannot be dropped. Without it one is looking at the image of
an infinite disjoint union, and an infinite union of closed sets need be
neither closed nor constructible; indeed the whole difficulty of the Hodge
conjecture in families is that $\Sigma$ itself is not known to be either. The
relative Chow variety of effective $n$-cycles of bounded degree would serve
equally well in place of the Hilbert schemes, with the same proof, and the
choice between them is immaterial here because all that is used is properness
over the base together with the local constancy of the cycle class. What is
used later, and what forces the bookkeeping with tuples rather than a single
subscheme, is that a rational cycle is a rational combination of several
integral ones and the several may have different Hilbert polynomials.} and
\[
  p^{-1}\Bigl(\,\bigcup_{\underline{P},\underline{q}}
    \Sigma_{\underline{P},\underline{q}}\Bigr)\;=\;\Sigma ,
\]
the union running over a countable index set.
\end{theorem}

\begin{proof}
Each $\Hilb^{P_{i}}(\cA/\cS)$ is projective over $\cS$ by Grothendieck's
existence theorem \cite{Gro61}, so the fibre product $\cH_{\underline{P}}$ is
projective over $\cS$, and being of finite type it has finitely many connected
components.

Fix a connected component $T$, replace it by its reduction, which changes
neither the underlying space nor any cohomology group, and let
$\cZ_{i}\subset\cA\times_{\cS}T$ be the pullback of the $i$-th universal
subscheme. It is flat over $T$ and its fibres have Hilbert polynomial $P_{i}$
of degree $n$, so the associated $n$-cycles, the top-dimensional components
counted with their multiplicities, are the fibres of a cycle on
$\cA\times_{\cS}T$ whose class restricts on the fibre over $t$ to $[Z_{i,t}]$.
Restriction of a global class to the fibres of a topologically
locally trivial family is a flat section of the local system, so
$t\mapsto\sum_{i}q_{i}[Z_{i,t}]$ is flat on $T$; and $\omega$ is flat on $\cS$
by \Cref{lem:omegadescends}, hence flat after pullback. Two flat sections of a
local system over a connected base agree everywhere as soon as they agree at
one point, so the locus where they agree is a union of connected components of
$\cH_{\underline{P}}$. That locus is therefore itself projective over $\cS$,
and the image of a proper morphism is Zariski closed.

For the equality of sets, let $s\in\cS$ with $\omega_{s}$ algebraic. Then
$\omega_{s}=\sum_{i}q_{i}[Z_{i}]$ for finitely many closed integral subschemes
$Z_{i}\subset A_{s}$ of dimension $n$ and rationals $q_{i}$; taking
$P_{i}$ to be the Hilbert polynomial of $Z_{i}$ puts $s$ in
$\Sigma_{\underline{P},\underline{q}}$. The converse holds because the class
of a closed subscheme of pure dimension $n$ is algebraic. Countability is
clear: a Hilbert polynomial is determined by its finitely many integer
coordinates in the binomial basis, and $\QQ$ is countable.
\end{proof}

\begin{corollary}[All, or almost nothing]\label{cor:allornothing}
Exactly one of the following holds.
\begin{enumerate}[label=\textup{(\alph*)},nosep]
\item $\Sigma=\cD_{n,n}$, and the Hodge conjecture holds in every degree for
  every abelian $2n$-fold of $(K,-1,n)$-Weil type with $\Hg(A)=\SU(V,H)$ in
  the family.
\item $\Sigma$ has empty interior in the analytic topology. Its complement is
  then a dense $G_{\delta}$, so $\omega_{s}$ fails to be algebraic for every
  $s$ outside a meagre set, and the Hodge conjecture fails.
\end{enumerate}
In particular, if $\omega_{s}$ is algebraic for all $s$ in some nonempty open
subset of $\cD_{n,n}$, then it is algebraic for all $s$.
\end{corollary}

\begin{proof}
Suppose $\Sigma$ contains a nonempty open set; then so does its image in
$\cS$, and we may pick a nonempty connected open $U\subseteq\cS$ contained in
that image. By \Cref{thm:closedlocus} the countably many sets
$\Sigma_{\underline{P},\underline{q}}\cap U$ are closed in $U$ and cover it.
A connected complex manifold is a Baire space, so one of them has nonempty
interior in $U$. A Zariski closed subset of the irreducible variety $\cS$
either is all of $\cS$ or is a proper closed subvariety, and a proper closed
subvariety of a connected complex manifold has empty interior. Hence that
$\Sigma_{\underline{P},\underline{q}}$ is all of $\cS$, so $\Sigma=\cD_{n,n}$
and (a) holds by \Cref{thm:residual}(a). If $\Sigma$ has empty interior, its
complement is the intersection of the countably many dense open sets
$\cD_{n,n}\setminus p^{-1}(\Sigma_{\underline{P},\underline{q}})$, and Baire
again gives that this intersection is dense.
\end{proof}

\begin{figure}[htbp]
\centering
\includegraphics[width=0.88\textwidth]{fig_properness}
\caption{The dichotomy of \Cref{cor:allornothing}. Each stratum
$\Sigma_{\underline{P},\underline{q}}$ is Zariski closed by
\Cref{thm:closedlocus}, and there are countably many of them; the base point
$s_{0}$ of \Cref{thm:everyfamily} lies on one. By \Cref{prop:densealg} the
union of the strata is dense, so a point $s$ off all of them is squeezed
between strata however small a neighbourhood is taken. The conjecture for the
family is the assertion that no such $s$ exists, and by
\Cref{cor:allornothing} it is enough that one stratum has nonempty interior.}
\label{fig:properness}
\end{figure}

The next fact says that case (b) of \Cref{cor:allornothing}, if it occurs,
occurs in an uncomfortable way: the algebraic locus is always dense, so the
two sets in the dichotomy would both be dense and the distinction between them
would be entirely one of measure rather than of position.

\begin{lemma}[Algebraicity is a rational invariant]\label{lem:Ginvariance}
$\Sigma$ is stable under the action of $G(\QQ)$ on $\cD_{n,n}$.
\end{lemma}

\begin{proof}
A point of $\cD_{n,n}$ is a complex structure $J$ on $V_{\RR}$ compatible with
$H$, and $g\in G(\QQ)$ carries $J$ to $gJg^{-1}$. The map $g$ is then
$\QQ$-rational and $\CC$-linear from $(V_{\RR},J)$ to $(V_{\RR},gJg^{-1})$, so
after clearing denominators it induces an isogeny $A_{J}\to A_{gJ}$, whose
graph is an algebraic correspondence. Algebraicity of a rational cohomology
class is preserved by such a correspondence and by its transpose. Finally $g$
lies in $\SU(V,H)$, so by the computation in \Cref{lem:omegadescends} it fixes
the Weil line pointwise, and the class transported from $\omega_{J}$ is
$\omega_{gJ}$.
\end{proof}

\begin{proposition}[The algebraic locus is dense]\label{prop:densealg}
Let $n\ge2$. Then $\Sigma$ is dense in $\cD_{n,n}$ in the analytic topology,
and its image in $\cS$ is dense.\footnote{We should be clear about which
of the two standard density statements this is, because they behave in
opposite ways. For a class that is not flat, the locus where it stays Hodge is
a countable union of proper closed subvarieties, and the content of
\Cref{thm:cdk} is that each of them is algebraic; density there would be a
statement about how special points distribute. For the class $\omega$ of
\Cref{setup:family} the Hodge locus is the whole family by
\Cref{prop:weilline}(ii), so \Cref{thm:cdk} says nothing at all, and the set
that carries information is the algebraic locus instead. The proposition below
is a statement about that set, and its proof uses no Hodge theory: the
algebraic locus is a union of orbits of a group whose rational points are
dense in its real points, and the density is inherited from the group.}
\end{proposition}

\begin{proof}
The hermitian form $H$ has rank $2n\ge4$ over $K$ and signature $(n,n)$, so it
is isotropic over $\RR$; over every $\QQ_{p}$ a hermitian form of rank at
least three over a quadratic extension is isotropic; and hermitian forms over
an imaginary quadratic field satisfy the Hasse principle
\cite[Ch.~10, \S 6]{Sch85}. So $H$ is isotropic over $\QQ$ and $G$ is
$\QQ$-isotropic. The group $G=\SU(V,H)$ is simply connected and almost
$\QQ$-simple, so by Platonov's theorem $G(\QQ)$ is generated by the rational
points of the unipotent radicals of its proper parabolic $\QQ$-subgroups
\cite[Thm.~7.6 and \S 7.2]{PR94}. For a unipotent group $U$ over $\QQ$ the
exponential identifies $U(\QQ)$ with a $\QQ$-vector space inside the
corresponding $\RR$-vector space $U(\RR)$, so $U(\QQ)$ is dense in $U(\RR)$.
The closure of $G(\QQ)$ in $G(\RR)$ therefore contains $U(\RR)$ for every such
$U$, and those generate the identity component $G(\RR)^{+}$, which acts
transitively on $\cD_{n,n}$. Hence every $G(\QQ)$-orbit in $\cD_{n,n}$ is
dense.

By \Cref{thm:everyfamily} the set $\Sigma$ is nonempty, and by
\Cref{lem:Ginvariance} it is $G(\QQ)$-stable, so it contains a dense orbit.
The projection $p$ is open and surjective, so the image is dense.
\end{proof}

\begin{corollary}[The conjecture for Weil type is a non-meagreness
statement]\label{cor:meagre}
Let $n\ge2$. The Hodge conjecture holds for every abelian $2n$-fold of
$(K,-1,n)$-Weil type in the family of \Cref{setup:family} if and only if
$\Sigma$ is non-meagre in $\cD_{n,n}$, and if and only if $\Sigma$ has
nonempty interior.
\end{corollary}

\begin{proof}
If $\Sigma=\cD_{n,n}$ both conditions hold. If not, \Cref{cor:allornothing}(b)
says $\Sigma$ is a countable union of closed sets with empty interior, hence
meagre, and \Cref{thm:residual}(b) says the conjecture fails.
\end{proof}

\subsection{The first criterion: a bound along one orbit}

The proof of \Cref{prop:densealg} produces the dense algebraic locus by moving
one cycle around by rational isogenies, and an isogeny multiplies degrees. So
the dense set it produces is stratified by unbounded data, which is exactly
what stops \Cref{cor:allornothing} from applying. Reversing this gives a
criterion, and it is the sharpest consequence of \Cref{thm:closedlocus} we
know how to state.

\begin{theorem}[A bound along the orbit closes the family]\label{thm:degreebound}
Let $n\ge2$ and let $s_{0}\in\Sigma$ be one of the base points supplied by
\Cref{thm:everyfamily}. Suppose that there is a finite set $\mathcal{T}$ of pairs
$(\underline{P},\underline{q})$ such that for every $g\in G(\QQ)$ the class
$\omega_{g\cdot s_{0}}$ is represented by some cycle whose Hilbert data lies
in $\mathcal{T}$. Then $\Sigma=\cD_{n,n}$, and the Hodge conjecture holds in every
degree for every abelian $2n$-fold of $(K,-1,n)$-Weil type with
$\Hg(A)=\SU(V,H)$ in that family.
\end{theorem}

\begin{remark}[The cycle is chosen, not transported]\label{rem:chosencycle}
The hypothesis quantifies over representing cycles at each point of the orbit
and asks only that one of them, at each point, be of bounded complexity. An
earlier form of this statement fixed the cycle, taking at $g\cdot s_{0}$ the
transport of the cycle at $s_{0}$ along the isogeny of
\Cref{lem:Ginvariance}. That form is stronger, and \Cref{thm:p1forces} below
shows that it is too strong to be usable: for a base cycle with a component of
finite stabiliser the transported data is unbounded, so its hypothesis is never
satisfied. The proof below uses only the weaker hypothesis, which is what the
covering argument needs.
\end{remark}

\begin{proof}
The orbit $G(\QQ)\cdot s_{0}$ is dense by the proof of \Cref{prop:densealg},
and by hypothesis it is the union of the finitely many sets
$G(\QQ)\cdot s_{0}\cap p^{-1}(\Sigma_{\underline{P},\underline{q}})$ with
$(\underline{P},\underline{q})\in\mathcal{T}$. A finite
union of closed sets is dense only if one of them is, so one
$p^{-1}(\Sigma_{\underline{P},\underline{q}})$ is dense; it is closed in the
analytic topology, hence equals $\cD_{n,n}$. Now apply
\Cref{thm:residual}(a).
\end{proof}

The hypothesis of \Cref{thm:degreebound} can be sharpened, and the sharpening
says which half of the Hilbert data is at issue. The leading coefficient is
not.

\begin{proposition}[The leading coefficient is constant on the orbit]
\label{prop:leadingconstant}
Let $s_{0}\in\Sigma$ and $g\in G(\QQ)$, and let $\omega_{s}$ denote the Weil
class at $s$ and $\eta$ the polarisation class. Then
\[
  \int_{A_{g\cdot s_{0}}}\omega_{g\cdot s_{0}}\cdot\frac{\eta^{n}}{n!}
  \;=\;\int_{A_{s_{0}}}\omega_{s_{0}}\cdot\frac{\eta^{n}}{n!},
\]
so the degree of the transported class with respect to the polarisation, which
is the leading coefficient of the Hilbert polynomial of any cycle
representing it, is the same at every point of the orbit.
\end{proposition}

\begin{proof}
By \Cref{lem:Ginvariance} the correspondence transporting the class is induced
by a $\QQ$-rational map which lies in $\SU(V,H)$, so it preserves the
hermitian form and hence the polarisation class, and by
\Cref{lem:omegadescends} it fixes the Weil line pointwise. An intersection
number of classes preserved by a correspondence is preserved by it. The
identification of the number with the leading coefficient is
Hirzebruch--Riemann--Roch on $A$, the higher coefficients involving the
structure sheaf of the cycle and not only its class.
\end{proof}

\begin{remark}[Where the unboundedness can only come from]
\label{rem:whereunbounded}
Write a representing cycle at a point of the orbit, the transported cycle
or any other, as a difference $\xi^{+}-\xi^{-}$ of effective cycles. By \Cref{prop:leadingconstant} the difference of their degrees is
constant along the orbit, and by the boundedness of the Chow variety of
effective cycles of bounded degree in a fixed polarised variety, if
$\deg\xi^{+}$ were bounded along the orbit then only finitely many tuples
$(\underline{P},\underline{q})$ could occur and \Cref{thm:degreebound} would
apply. So the hypothesis of that theorem is equivalent to a bound on the
effective part alone.

That is not a formality, and the mechanism that threatens it is visible. The
transport at $g$ is by the isogeny obtained from $g$ by clearing denominators;
if $g$ has denominator $m$ on the reference lattice, which has rank $4n$, the
isogeny has degree a divisor of $m^{4n}$, and pushing a cycle forward along it
multiplies the multiplicities by factors of that order. The rational points of
$G$ are dense precisely because their denominators are unbounded, so along any
orbit dense enough to be useful the isogeny degrees are unbounded as well.
What \Cref{thm:degreebound} asks is that the cycles nevertheless simplify: that
at each point of the orbit the class $\omega$, whose degree has not changed,
admits an effective representation no more complicated than before. For the
transported cycle itself the answer is negative (\Cref{thm:p1forces} below),
and for a cycle chosen at each point the bound is equivalent to the conclusion
(\Cref{thm:p1equivalent}), so it cannot be decided before the conclusion is.
The content of the present remark is only that the question has been reduced from the whole
Hilbert datum to the degree of the effective part, and that the part of the
datum which is determined by cohomology is already constant.
\end{remark}

\subsection{What the transport does, exactly}

The mechanism of \Cref{rem:whereunbounded} can be computed rather than
described, and computing it decides one of the two readings of (P1). Write
$\Lambda$ for the reference lattice, of rank $4n$ over $\ZZ$, and $E$ for the
alternating form of the polarisation, so that a point $s$ of $\cD_{n,n}$ is a
complex structure on $\Lambda\otimes\RR$ and $A_{s}=(\Lambda\otimes\RR)/\Lambda$.
Throughout, $G=\dim A_{s}=2n$ and a component of a cycle representing
$\omega_{s}$ has dimension $n$.

\begin{lemma}[The transport, in coordinates]\label{lem:transport}
Let $g\in G(\QQ)$ and let $c=c(g)$ be the least positive integer with
$cg\Lambda\subseteq\Lambda$. Then $\varphi_{g}:=cg$ induces an isogeny
$A_{s_{0}}\to A_{g\cdot s_{0}}$ with
\begin{enumerate}[label=\textup{(\roman*)},nosep,leftmargin=2.6em]
\item $\deg\varphi_{g}=c^{4n}$;
\item $\varphi_{g}^{*}\eta_{g\cdot s_{0}}=c^{2}\eta_{s_{0}}$ and
  $\varphi_{g}^{*}\omega_{g\cdot s_{0}}=c^{2n}\omega_{s_{0}}$;
\item the transported cycle representing $\omega_{g\cdot s_{0}}$ is
  $c^{2n-4n}\,\varphi_{g*}(Z)$, where $Z$ represents $\omega_{s_{0}}$.
\end{enumerate}
Moreover $\sup_{g\in G(\QQ)}c(g)=\infty$.
\end{lemma}

\begin{proof}
(i) The index of $cg\Lambda$ in $\Lambda$ is $|\det_{\QQ}(cg)|
=c^{4n}|\det_{\QQ}g|$, and $\det_{\QQ}g=\Nm_{K/\QQ}(\det_{K}g)=1$ for
$g\in\SU(V,H)$.

(ii) $g$ preserves $E$, so $(cg)^{*}E=c^{2}E$, which is the first identity.
On $H^{2n}=\bigwedge^{2n}(\Lambda^{\vee}\otimes\QQ)$ the map
$\varphi_{g}^{*}$ is $c^{2n}\bigwedge^{2n}g^{\vee}$, and $g\in\SU(V,H)$ fixes
the Weil line pointwise by \Cref{lem:omegadescends}.

(iii) $\varphi_{g*}\varphi_{g}^{*}=\deg\varphi_{g}$, so
$\omega_{g\cdot s_{0}}=(\deg\varphi_{g})^{-1}\varphi_{g*}
\varphi_{g}^{*}\omega_{g\cdot s_{0}}
=c^{2n-4n}\varphi_{g*}\omega_{s_{0}}$ by (i) and (ii).

For the last sentence, $G$ is $\QQ$-isotropic by the proof of
\Cref{prop:densealg}, so it has a proper parabolic $\QQ$-subgroup whose
unipotent radical $U$ satisfies $U(\QQ)\cong\QQ^{k}$ with $k\ge1$ through the
exponential; the entries of $u\in U(\QQ)$ are $\QQ$-linear in that coordinate,
so $c(u)$ is unbounded on $U(\QQ)$. Item (XXXVII) of \Cref{sec:verification}
verifies (i) and (ii) on an explicit sample of such $u$ with $c$ up to $29$,
for $2n=4,6,8$.
\end{proof}

\begin{lemma}[The multiplicity is the stabiliser met by the kernel]
\label{lem:multiplicity}
Let $\varphi\colon A\to A'$ be an isogeny of abelian varieties and $Z\subseteq A$
an integral closed subvariety. Put
$G(Z)=\{x\in A:Z+x=Z\}$, a closed algebraic subgroup. Then
\[
  \varphi_{*}[Z]=m\,[\varphi(Z)],\qquad
  m=\bigl|\ker\varphi\cap G(Z)\bigr| .
\]
If moreover $\dim Z=n$ and $\varphi^{*}\eta'=c^{2}\eta$, then
\[
  \deg_{\eta'}\varphi(Z)=\frac{c^{2n}}{m}\,\deg_{\eta}Z .
\]
\end{lemma}

\begin{proof}
The map $\varphi|_{Z}\colon Z\to\varphi(Z)$ is finite and surjective, and its
degree $m$ is the number of points of a general fibre. For $z\in Z$ the fibre
through $z$ is $Z\cap(z+\ker\varphi)$, which as a set of translations is
$\{t\in\ker\varphi:z+t\in Z\}$. For a fixed $t$ the locus
$\{z\in Z:z+t\in Z\}=Z\cap(Z-t)$ is all of $Z$ when $t\in G(Z)$ and a proper
closed subset otherwise; since $\ker\varphi$ is finite, a general $z$ avoids
the finitely many proper subsets, and there
$\{t\in\ker\varphi:z+t\in Z\}=\ker\varphi\cap G(Z)$. This is the first
identity, and $\varphi_{*}[Z]=m[\varphi(Z)]$ by the definition of the
pushforward of cycles.

For the degree, the projection formula gives
\[
  m\deg_{\eta'}\varphi(Z)=\int\varphi_{*}[Z]\cdot\frac{\eta'^{n}}{n!}
  =\int[Z]\cdot\frac{(\varphi^{*}\eta')^{n}}{n!}
  =c^{2n}\int[Z]\cdot\frac{\eta^{n}}{n!}=c^{2n}\deg_{\eta}Z .
\]
Item (XXXVII) verifies both identities for subtori of an abelian fourfold,
with $m$ computed as a lattice index by Smith normal form and the degrees as
Pfaffians; there the two cases are visible side by side, a subtorus preserved
by $\varphi$ having $m=c^{2n}$ and image degree constant, and one not
preserved having $m=c^{2n-1}$ and image degree growing linearly in $c$.
\end{proof}

\begin{theorem}[What \textup{(P1)} forces on the base cycle]\label{thm:p1forces}
Let $s_{0}\in\Sigma$ and let $Z=\sum_{i}q_{i}Z_{i}$, with $q_{i}\ne0$ and the
$Z_{i}$ integral of dimension $n$, represent $\omega_{s_{0}}$.
\begin{enumerate}[label=\textup{(\roman*)},leftmargin=2.6em,itemsep=2pt]
\item For $g\in G(\QQ)$ the transported cycle is
  $\sum_{i}c(g)^{2n-4n}m_{i}(g)\,q_{i}\,[\varphi_{g}(Z_{i})]$ with
  $m_{i}(g)=|\ker\varphi_{g}\cap G(Z_{i})|$, and the $i$-th component has
  degree $c(g)^{2n}\deg_{\eta}(Z_{i})/m_{i}(g)$.
\item If some $Z_{i}$ has finite stabiliser $G(Z_{i})$, then
  $\deg\varphi_{g}(Z_{i})\ge c(g)^{2n}\deg_{\eta}(Z_{i})/|G(Z_{i})|$, which is
  unbounded on $G(\QQ)$. The Hilbert data of the transported cycle then takes
  infinitely many values, and the hypothesis of \Cref{thm:degreebound} in the
  form that fixes the transported cycle is not satisfied by $Z$.
\item Consequently, a base cycle for which the transported data is bounded
  must have every component of positive-dimensional stabiliser, that is, every
  $Z_{i}$ must be stable under translation by a positive-dimensional abelian
  subvariety of $A_{s_{0}}$.
\end{enumerate}
\end{theorem}

\begin{proof}
(i) is \Cref{lem:transport}(iii) and \Cref{lem:multiplicity}.

(ii) $\ker\varphi_{g}\cap G(Z_{i})\subseteq G(Z_{i})$, so
$m_{i}(g)\le|G(Z_{i})|$, a constant; the displayed bound is (i), and $c(g)$ is
unbounded by \Cref{lem:transport}. Distinct degrees of the components give
distinct Hilbert polynomials, so the data takes infinitely many values.

(iii) is the contrapositive of (ii), together with the fact that an algebraic
subgroup of positive dimension in an abelian variety has an abelian subvariety
of positive dimension as its identity component.
\end{proof}

\begin{remark}[What this settles and what it leaves]\label{rem:p1meaning}
Three things follow, and we separate them.

The statement (P1) has two readings, and one of them is now closed. Read with
the cycle fixed at $s_{0}$ and transported, it fails for every base cycle this
paper writes down: those cycles are intersections of divisors, and an
intersection of $n$ divisors on an abelian $2n$-fold has finite stabiliser
unless the divisors are pulled back from a quotient, which the ones here are
not. Read with the cycle chosen afresh at each point of the orbit, which is
the form \Cref{thm:degreebound} now takes, it is untouched by
\Cref{thm:p1forces} and remains open.

The two readings are different, and the difference is the content of
\Cref{lem:multiplicity}. The class $\omega_{g\cdot s_{0}}$ has constant degree
along the orbit, by \Cref{prop:leadingconstant}; what grows is the complexity
of the particular representative the transport produces, and it grows because
the isogeny spreads each component out by the factor $c^{2n}/m_{i}$. A
representative of bounded complexity may well exist at every point; the
transport does not find it.

The third is a demand on any future attempt at (P1). It must either produce a
base cycle supported on subvarieties with positive-dimensional stabiliser,
which for a Weil class is a strong and possibly impossible requirement, or
abandon the transport and bound instead the minimal complexity of a
representative, which is a statement about the Chow variety of the family and
not about any one cycle. We do not know how to do either.
\end{remark}

The bound on the minimal complexity is the form in
which \Cref{thm:degreebound} now stands. It is the weakest hypothesis under
which the covering argument runs, and it is therefore the form in which (P1)
would have to be attacked. The next theorem decides it, and the answer is that
there is nothing there to attack: the hypothesis is not a sufficient condition
for the conjecture on the family but another way of writing it.

\begin{theorem}[The bounded form is the conjecture itself]\label{thm:p1equivalent}
Let $n\ge2$, let $s_{0}\in\Sigma$ and keep \Cref{setup:family}. The following
are equivalent.
\begin{enumerate}[label=\textup{(\alph*)},leftmargin=2.6em,itemsep=2pt]
\item There is a finite set $\mathcal{T}$ of pairs
  $(\underline{P},\underline{q})$ such that at every point of the orbit
  $G(\QQ)\cdot s_{0}$ the class $\omega$ is represented by a cycle whose
  Hilbert data lies in $\mathcal{T}$; that is, \textup{(P1)} holds at
  $s_{0}$.
\item There is a single pair $(\underline{P},\underline{q})$ such that at
  every point of $\cD_{n,n}$, not merely of the orbit, the class $\omega$ is
  represented by a cycle with that datum.
\item $\Sigma=\cD_{n,n}$.
\item $\mathbf{W}(K,n,\delta)$ holds, that is, the Weil classes are algebraic
  on every member of the family.
\end{enumerate}
In particular \Cref{thm:degreebound} is an equivalence, and \textup{(P1)} is a
reformulation of the conjecture for the family and not a reduction of it. The
same holds for every CM field, with \Cref{thm:cmpropagation}(iv) in place of
\Cref{thm:degreebound}.
\end{theorem}

\begin{proof}
(a) implies (c) is \Cref{thm:degreebound}. (b) implies (a) is immediate, the
orbit being contained in $\cD_{n,n}$ and $\mathcal{T}$ being the one pair.
(c) is equivalent to (d) by \Cref{thm:residual}(a) and (b).

It remains to deduce (b) from (c). If $\Sigma=\cD_{n,n}$ then $\Sigma$ has
nonempty interior, and the first paragraph of the proof of
\Cref{cor:allornothing} shows that some
$\Sigma_{\underline{P},\underline{q}}$ is then all of $\cS$: the countably
many sets $\Sigma_{\underline{P},\underline{q}}\cap U$ are closed in a
nonempty connected open $U\subseteq\cS$ and cover it, a connected complex
manifold is a Baire space, so one of them has nonempty interior in $U$, and a
proper Zariski closed subvariety of the irreducible variety $\cS$ has empty
interior. By the definition of $\Sigma_{\underline{P},\underline{q}}$ in
\Cref{thm:closedlocus} as the image of a proper morphism, every $s\in\cS$
then carries a tuple $\underline{Z}$ of subschemes with Hilbert polynomials
$\underline{P}$ and $\sum_{i}q_{i}[Z_{i}]=\omega$, which is (b).

For a CM field the same argument applies verbatim, using
\Cref{thm:cmpropagation}(iii) in place of \Cref{thm:closedlocus} and
\Cref{cor:allornothing}.
\end{proof}

\begin{remark}[What is left of the first criterion]\label{rem:p1collapse}
\Cref{thm:p1equivalent} removes (P1) from the list of things that would close
the family, and we should be exact about what it does and does not say.

It does not say that \Cref{thm:degreebound} is false or that its proof is
defective. The implication it proves is correct, and the covering argument is
the only place in this paper where density of the rational orbit is converted
into a global statement. What the theorem says is that the implication has a
converse, so that the hypothesis is no weaker than the conclusion. A criterion
of that kind cannot be used: verifying it at one family is exactly as hard as
proving the conjecture for that family, no more and no less.

Nor does it say that the boundedness is uninteresting. It is a concrete
finiteness statement about the Chow variety of the universal family, and it is
the only reformulation of $\mathbf{W}(K,n,\delta)$ we know that mentions
neither Hodge theory nor deformation theory. What it is not is a step.

The two forms of the propagation statement are therefore not on the same
footing, and we had them on the same footing until now. (P2) remains a genuine
sufficient condition: injectivity of one semiregularity map at one point
implies the conjecture for the whole family, and nothing here shows the
converse, nor do we see how it could, since $\Sigma=\cD_{n,n}$ says nothing
about the deformation theory of any particular perfect complex. So the
propagation half of the frontier is the single statement (P2), and the
constraints \Cref{thm:p2forces} and \Cref{cor:p2shape} place on it are the
constraints on the whole of that half.
\end{remark}

At the quaternionic points of \Cref{setup:quat} the mechanism of
\Cref{rem:whereunbounded} can be seen explicitly. Take $n=2$ and the integral
lattice
$\Lambda_{b}=\cO_{b}^{2}$, $\cO_{b}=\ZZ\langle 1,i,j,ij\rangle$, for
$b\in\ZZ_{>0}$. Put
\[
  g_{b}(x,y)=E_{b}(jx,y),\qquad
  \lambda_{b}(x,y)=-\tfrac{1}{d}\,g_{b}(ix,y),\qquad
  Z_{b}=g_{b}^{2}-d\lambda_{b}^{2},\qquad Z'_{b}=2g_{b}\lambda_{b},
\]
so that $Z_{b}+\sqrt{-d}\,Z'_{b}=(g_{b}+\sqrt{-d}\,\lambda_{b})^{2}$ and
$Z_{b},Z'_{b}$ span the Weil plane. Write $x\in\cO_{b}$ as $x=u+vj$ with
$u,v\in\ZZ[i]$; the pair $(\Lambda_{b},i)$ is then the same $\ZZ[i]$-module
$\ZZ[i]^{4}$ for every $b$, and we use this identification throughout.

\begin{proposition}[The quaternionic Weil cycle is a square multiple]
\label{prop:quatdivisible}
With these data, for every $b\in\ZZ_{>0}$:
\begin{enumerate}[label=\textup{(\roman*)},nosep]
\item $E_{b}=E_{u}+b\,E_{v}$, $g_{b}=b\,g_{1}$ and $\lambda_{b}=b\,\lambda_{1}$,
  where $E_{u},E_{v},g_{1},\lambda_{1}$ are integral on $\ZZ[i]^{4}$ and do not
  depend on $b$; in particular $g_{1},\lambda_{1}\in\mathrm{NS}(A_{b})$;
\item $Z_{b}=b^{2}Z_{1}$ and $Z'_{b}=b^{2}Z'_{1}$ in $\mathrm{CH}^{2}(A_{b})$
  modulo algebraically trivial cycles, for every choice of line bundles
  representing $g_{b},\lambda_{b}$; in cohomology the equalities hold in
  $\wedge^{4}\Lambda^{*}$;
\item $\int\eta_{b}^{4}=b^{2}\int\eta_{1}^{4}$,
  $\int g_{b}^{2}\eta_{b}^{2}=-\tfrac{b}{3}\int\eta_{b}^{4}$ and
  $\int Z_{b}^{2}=\tfrac{4b^{2}}{3}\int\eta_{b}^{4}$, so that
  $(\int g_{b}^{2}\eta_{b}^{2})^{2}=\tfrac{1}{12}\int Z_{b}^{2}\int\eta_{b}^{4}$;
\item for $\alpha\in K^{\times}$ the element $\alpha j$ generates the same
  algebra at the same point, and replacing $j$ by $\alpha j$ multiplies the
  two ratios in \textup{(iii)} by $N(\alpha)$ and $N(\alpha)^{2}$.
\end{enumerate}
\end{proposition}

\begin{proof}
For $x=u+vj$ one has $\bbar{x}=\bbar{u}-vj$, and the relations $ju=\bbar{u}j$,
$j^{2}=b$ give
$\bbar{x}\,i\,y\equiv\bbar{u}\,i\,u'+b\,v\,i\,\bbar{v'}$ modulo $Kj$. Since
$\mathrm{trd}$ vanishes on $Kj$ and restricts to $\mathrm{tr}_{K/\QQ}$ on $K$,
\[
  E_{b}(x,y)=\textstyle\sum_{k}a_{k}\bigl(\mathrm{tr}(i\,\bbar{u_{k}}u'_{k})
  +b\,\mathrm{tr}(i\,v_{k}\bbar{v'_{k}})\bigr).
\]
From $jx=b\bbar{v}+\bbar{u}j$ one reads off
$g_{b}(x,y)=b\sum_{k}a_{k}\bigl(\mathrm{tr}(i\,v_{k}u'_{k})
+\mathrm{tr}(i\,\bbar{u_{k}}\bbar{v'_{k}})\bigr)$, which is $b\,g_{1}$, and
$\lambda_{b}=b\lambda_{1}$ follows. The traces take integral values on
$\ZZ[i]$, and $g_{1}=b^{-1}g_{b}$ is a rational multiple of a Hodge class,
which proves (i). Then $g_{b}+\sqrt{-d}\,\lambda_{b}=b(g_{1}+\sqrt{-d}\,\lambda_{1})$
and (ii) holds in cohomology. Choosing $L,M$ with $c_{1}(L)=g_{1}$,
$c_{1}(M)=\lambda_{1}$, the cycle $c_{1}(L^{b})^{2}-d\,c_{1}(M^{b})^{2}$ equals
$b^{2}(c_{1}(L)^{2}-d\,c_{1}(M)^{2})$ in $\mathrm{CH}^{2}(A_{b})$; another
choice changes both sides by algebraically trivial cycles.

For (iii), $E_{u}$ and $E_{v}$ live on complementary halves $U,V'$ of
$\Lambda\otimes\QQ$ of rank four, so $\eta_{u}^{3}=\eta_{v}^{3}=0$ and
$\int\eta_{b}^{4}=6b^{2}\int\eta_{u}^{2}\eta_{v}^{2}$. The form $g_{1}$ pairs
$U$ with $V'$, so $g_{1}^{2}\in\wedge^{2}U^{*}\otimes\wedge^{2}V'^{*}$, and
only the mixed term survives in
$\int g_{b}^{2}\eta_{b}^{2}=2b^{3}\int g_{1}^{2}\eta_{u}\eta_{v}$; likewise
$\int Z_{b}^{2}=b^{4}\int Z_{1}^{2}$. Hence the ratios
$\int g_{b}^{2}\eta_{b}^{2}/\int\eta_{b}^{4}$ and
$\int Z_{b}^{2}/\int\eta_{b}^{4}$ are $b$ and $b^{2}$ times their values at
$b=1$, which are $-1/3$ and $4/3$ by direct evaluation. For (iv), $(\alpha j)^{2}=N(\alpha)b$ and
$g_{\alpha j}+\sqrt{-d}\,\lambda_{\alpha j}$ is $\alpha$ or $\bbar{\alpha}$
times $g_{j}+\sqrt{-d}\,\lambda_{j}$.
\end{proof}

Item (XXIII) of \Cref{sec:verification} confirms (i) to (iv) in exact
arithmetic for $d\in\{1,3\}$, two weight vectors and fifteen values of $b$ up
to $101$. It also computes the integral Weil lattice
$\HW_{\ZZ}=(\QQ Z_{b}+\QQ Z'_{b})\cap\wedge^{4}\Lambda^{*}$ and the
sublattice spanned inside it by products of two integral divisor classes: for
$d\in\{1,2,3,7\}$, three weight vectors and seven values of $b$, the Gram
matrix of $\HW_{\ZZ}$ is $\mathrm{diag}(8d,8d^{2})$ and the index is
$2(a_{1}a_{2})^{2}$, both independent of $b$.

\begin{remark}[The multiple sits in the polarisation]\label{rem:quatdivisible}
At these points the Weil class carries no multiple that grows with $b$ once
it is written in integral terms: by \Cref{prop:quatdivisible}(ii) it is
$b^{2}Z_{1}$, with $Z_{1}$ one integral algebraic cycle built from the two
fixed divisor classes $g_{1},\lambda_{1}$. The factor that the orbit argument
would otherwise have to absorb into the coefficients $\underline{q}$ of
\Cref{thm:degreebound} is therefore absent. What grows is the polarisation:
$\int\eta_{b}^{4}=b^{2}\int\eta_{1}^{4}$, so the points $A_{b}$ carry
polarisations of unbounded degree, and an isogeny carrying $\eta_{b}$ to the
polarisation of one fixed family has degree $\int\eta_{b}^{4}/\int\eta^{4}$,
which grows like $b^{2}$. This is the unbounded isogeny degree of
\Cref{rem:whereunbounded}, located at an explicit family of points, and it is
untouched by the integrality of $Z_{1}$.
\end{remark}

The base points of \Cref{thm:everyfamily}, and every algebraicity statement
for Weil classes proved in this paper, represent $\omega$ through products of
divisor classes. The next result says that a representation of that kind can
hold with bounded degrees only on a nowhere dense set, so the bound asked for
in \Cref{thm:degreebound} has to come from cycles that are not products of
divisors. For a class $x\in H^{2}(A_{s},\RR)$ put
\[
  t(x)=\frac{\int x\,\eta^{2n-1}}{\int\eta^{2n}},\qquad
  I(x)=t(x)^{2}\!\int\eta^{2n}-\int\bigl(x-t(x)\eta\bigr)^{2}\eta^{2n-2}.
\]
Since $\cD_{n,n}$ is contractible, the local system $R^{2}\pi_{*}\ZZ$ is
trivial; write $H^{2}_{\ZZ}$ for its fibre, identified with $H^{2}(A_{s},\ZZ)$
at every $s$. Both $t$ and $I$ are polynomial functions on $H^{2}_{\ZZ}$ and
do not depend on $s$.

\begin{lemma}[The Hodge norm of a divisor]\label{lem:nefnorm}
Let $h_{s}$ be the Hodge norm on $H^{2}(A_{s},\RR)$ defined by the
polarisation, normalised by $h_{s}(\eta)=\int\eta^{2n}$. Then
\begin{enumerate}[label=\textup{(\roman*)},nosep]
\item $h_{s}(x)=I(x)$ for every real class $x$ of type $(1,1)$ at $s$;
\item if $D$ is an effective divisor on $A_{s}$ then
  $I([D])\le 2\bigl(\int[D]\eta^{2n-1}\bigr)^{2}\big/\!\int\eta^{2n}$.
\end{enumerate}
\end{lemma}

\begin{proof}
Write $x=t\eta+x_{0}$ with $t=t(x)$, so that $x_{0}$ is primitive. The
Lefschetz decomposition is orthogonal for $h_{s}$, and on primitive classes
$h_{s}(x_{0})=-\int x_{0}\wedge C_{s}x_{0}\wedge\eta^{2n-2}$, where $C_{s}$ is
the Weil operator, positive by the Hodge--Riemann relations. For $x_{0}$ of type
$(1,1)$, $C_{s}x_{0}=x_{0}$, which gives (i). An effective divisor on an
abelian variety is nef, since a translate of it meets any given curve
properly, so $\int[D]^{2}\eta^{2n-2}\ge0$. Expanding,
$\int[D]^{2}\eta^{2n-2}=t^{2}\int\eta^{2n}+\int x_{0}^{2}\eta^{2n-2}$, hence
$-\int x_{0}^{2}\eta^{2n-2}\le t^{2}\int\eta^{2n}$ and
$I([D])\le2t^{2}\int\eta^{2n}$, which is (ii).
\end{proof}

\begin{proposition}[Divisor products of bounded degree]\label{prop:divroute}
In \Cref{setup:family}, fix $C>0$ and let $\Sigma_{C}\subset\cD_{n,n}$ be the
set of points $s$ at which $\omega_{s}$ is a $\QQ$-linear combination of
classes $[D_{1}]\cdots[D_{n}]$ with every $D_{i}$ an effective divisor on
$A_{s}$ of degree $\int[D_{i}]\eta^{2n-1}\le C$. Then $\Sigma_{C}$ is
contained in a locally finite union of proper closed analytic subsets of
$\cD_{n,n}$. In particular $\Sigma_{C}$ is nowhere dense, and on any dense
subset of $\cD_{n,n}$ at which $\omega$ is a combination of products of
effective divisors the degrees of those divisors are unbounded on every
nonempty open set.
\end{proposition}

\begin{proof}
For $x\in H^{2}_{\ZZ}$ let $\mathrm{Hdg}(x)\subset\cD_{n,n}$ be the
closed analytic set where $x$ is of type $(1,1)$. If $x\notin\QQ\eta$ it is
proper, because $\NS(A_{s})\otimes\QQ=\QQ\eta$ at a very general point by
\Cref{lem:genericNS}. Let $s\in\Sigma_{C}$. Not every $[D_{i}]$ in the
representation lies in $\QQ\eta$, since $\QQ\eta^{n}$ and the Weil line carry
different characters by \Cref{lem:charsdiffer}. So $s\in\mathrm{Hdg}(x)$ for
some $x$ in
\[
  X_{C}=\bigl\{x\in H^{2}_{\ZZ}\setminus\QQ\eta\;:\;
  I(x)\le 2C^{2}/\textstyle\int\eta^{2n}\bigr\},
\]
by \Cref{lem:nefnorm}(ii). Now let $U\subset\cD_{n,n}$ be open with compact
closure and $s_{0}\in U$. The Hodge norm depends continuously on the point, so
$h_{s_{0}}\le c_{U}\,h_{s}$ for some $c_{U}$ and all $s\in U$. If
$\mathrm{Hdg}(x)$ meets $U$ at $s$ and $x\in X_{C}$, then by
\Cref{lem:nefnorm}(i)
$h_{s_{0}}(x)\le c_{U}h_{s}(x)=c_{U}I(x)\le 2c_{U}C^{2}/\int\eta^{2n}$, and
only finitely many integral classes satisfy this. So the sets
$\mathrm{Hdg}(x)$, $x\in X_{C}$, form a locally finite family of proper closed
analytic subsets. Their union is closed with empty interior in the connected
manifold $\cD_{n,n}$, and it contains $\Sigma_{C}$.
\end{proof}

\begin{remark}[What the proposition closes]\label{rem:divroute}
The loci on which this paper proves $\omega$ algebraic are the split,
quaternionic and product loci and their translates under $G(\QQ)$, and on
each of them $\omega$ is a rational polynomial in divisor classes, since an
isogeny carries $\NS\otimes\QQ$ isomorphically. \Cref{prop:divroute} shows that no
bounded family of such cycles covers a dense set, so the finiteness in
\Cref{thm:degreebound} cannot be obtained by bounding the divisors that
produce $\omega$ at the points where it is known to be algebraic. It has to
come from the transported cycles of \Cref{lem:Ginvariance} themselves, which
are push-forwards along isogenies and not products of divisors on the target.
That is the unbounded isogeny degree of \Cref{rem:whereunbounded} and
\Cref{rem:quatdivisible}, now shown to be unavoidable for the divisorial
constructions. The argument uses nothing about $\omega$ beyond the fact that it
is not a multiple of $\eta^{n}$, and nothing about the family beyond the
generic Picard number, so it applies equally to any variational argument that
feeds on divisor products.
\end{remark}

At $n=2$ the phenomenon can be exhibited on one fixed polarised lattice. Take
$\Lambda=\cO^{2}$ with $\cO=\ZZ\langle1,i,j,ij\rangle$, $i^{2}=-d$, $j^{2}=1$,
and $E$ as in \eqref{eq:quatE}; this is one member of the family, with its
lattice and polarisation fixed. Let
$R=\{x:\ x(iu,v)=x(u,iv)\}\subset H^{2}(\Lambda,\QQ)$ be the space of
$K$-bilinear classes, where the divisor classes of the quaternionic points
live, and for $x\in R$ let $\phi_{x}=E^{-1}x$, so that $x(u,v)=E(\phi_{x}u,v)$.

\begin{proposition}[A pencil of quaternionic loci]\label{prop:pencil}
In this setting, for $d\in\{1,3\}$:
\begin{enumerate}[label=\textup{(\roman*)},nosep]
\item $R\cap H^{2}(\Lambda,\ZZ)$ has rank $12$, every $x\in R$ has $t(x)=0$,
  and $I(x)=\tfrac{1}{24}\,\mathrm{tr}(\phi_{x}^{2})\int\eta^{4}$; the form $I$
  has signature $(8,4)$ on $R$;
\item there are $x,y\in R\cap H^{2}(\Lambda,\ZZ)$ such that
  $\phi_{x+ky}^{2}=\beta(k)$ is a positive rational scalar for every
  $k\in\ZZ$, with $\beta(k)=\tfrac12+\tfrac12k+\tfrac14k^{2}$ for $d=1$ and
  $\beta(k)=\tfrac19+\tfrac29k^{2}$ for $d=3$;
\item for every $k$ the set $\mathrm{Hdg}(x+ky)$ is a quaternionic locus
  with algebra $(-d,\beta(k))$, nonempty of dimension $3$, and loci with
  different values of $\beta$ are not related by any automorphism of
  $(\Lambda,E,i)$;
\item at a very general point of $\mathrm{Hdg}(x+ky)$ one has
  $\NS\otimes\QQ=\QQ\eta\oplus\QQ x_{k}\oplus\QQ\,x_{k}(i\cdot,\cdot)$ with
  $x_{k}=x+ky$, and for every $k\in\ZZ$ the minimum $\mu_{k}$ of $I$ on
  $\NS\setminus\QQ\eta$ satisfies
  $\mu_{k}\ge\beta(k)\big/\bigl(3N_{d}^{2}\int\eta^{4}\bigr)$, with $N_{1}=2$
  and $N_{3}=4$.
\end{enumerate}
Consequently, at a very general point of $\mathrm{Hdg}(x+ky)$ every effective
divisor whose class is not a multiple of $\eta$ has degree at least
$\sqrt{\beta(k)/6}\,/N_{d}$, which grows linearly in $|k|$.
\end{proposition}

\begin{proof}
(i) and (ii) are finite exact computations, item (XXIV) of
\Cref{sec:verification}: the trace identity is checked on all pairs of a basis
of $R\cap H^{2}(\Lambda,\ZZ)$, which proves it on $R$ by bilinearity, and in
(ii) the three matrices $\phi_{x}^{2}$, $\phi_{y}^{2}$ and
$\phi_{x}\phi_{y}+\phi_{y}\phi_{x}$ are scalar, so
$\phi_{x+ky}^{2}=\phi_{x}^{2}+k(\phi_{x}\phi_{y}+\phi_{y}\phi_{x})+k^{2}\phi_{y}^{2}$
is scalar for every $k$; $\beta$ has positive leading coefficient and negative
discriminant. For (iii) write $\phi=\phi_{x+ky}$ and $\beta=\beta(k)$. The map
$\phi$ is $E$-self-adjoint, anticommutes with $i$, and $\phi^{2}=\beta>0$, so
$V_{\RR}=V_{+}\oplus V_{-}$ with $V_{\pm}$ the $\pm\sqrt{\beta}$ eigenspaces,
$E$-orthogonal to each other, and $i$ exchanges them. A complex structure $J$
on $V_{+}$ compatible with $E|_{V_{+}}$, extended to $V_{-}=iV_{+}$ by
$J(iw)=iJw$, commutes with $i$ and $\phi$ and is compatible with $E$, and $i$
has eigenvalues $\pm\sqrt{-d}$ each with multiplicity $2$ on $V^{1,0}$; so
these $J$ are points of $\cD_{2,2}$ at which $x+ky$ is of type $(1,1)$, and
they form a copy of the Siegel space of $V_{+}$, of dimension $3$, in
agreement with \Cref{prop:lagrangian}. An automorphism $\gamma$
of $(\Lambda,E,i)$ sends $\phi_{x}$ to $\gamma\phi_{x}\gamma^{-1}$ and so
preserves $\mathrm{tr}(\phi_{x}^{2})$.

For the first half of (iv), let $S_{k}\subset\cD_{2,2}$ be the Siegel space
just described and $\mathcal{C}_{k}\subset\Sp(V_{\RR},E)$ the centraliser of
$i$ and $\phi$. An element of $\mathcal{C}_{k}$ preserves $V_{+}$ and is
determined by its restriction there, because it commutes with $i$, so
$\mathcal{C}_{k}\cong\Sp(V_{+})$, acting on $V_{\RR}=V_{+}\oplus iV_{+}$ as
$g\oplus igi^{-1}$. A real class is of type $(1,1)$ at $J$ exactly when it is
invariant under $J$. The points of $S_{k}$ form one orbit
$\{gJ_{0}g^{-1}:g\in\mathcal{C}_{k}\}$ of a fixed point $J_{0}$. Let $z$ be a
real $2$-form invariant under every point of $S_{k}$, and write
$\mathfrak{c}=\mathrm{Lie}(\mathcal{C}_{k})\cong\mathfrak{sp}(V_{+})$,
$\mathfrak{k}$ for the centraliser of $J_{0}$ in $\mathfrak{c}$ and
$\mathfrak{p}=[\mathfrak{c},J_{0}]$. For $X\in\mathfrak{c}$ the curve
$g_{t}=\exp(tX)$ gives $(g_{t}J_{0}g_{t}^{-1})^{*}z=z$ for all $t$, and the
derivative at $t=0$ says that $[X,J_{0}]$ annihilates $z$, so $\mathfrak{p}$
annihilates $z$. The annihilator of $z$ in $\mathfrak{c}$ is a Lie
subalgebra, so it contains $\mathfrak{p}+[\mathfrak{p},\mathfrak{p}]$. Now
$\mathfrak{c}=\mathfrak{k}\oplus\mathfrak{p}$ with
$[\mathfrak{k},\mathfrak{p}]\subseteq\mathfrak{p}$ and
$[\mathfrak{p},\mathfrak{p}]\subseteq\mathfrak{k}$, so
$\mathfrak{p}+[\mathfrak{p},\mathfrak{p}]$ is an ideal of $\mathfrak{c}$;
it is nonzero, and $\mathfrak{sp}_{4}(\RR)$ is simple, so it is all of
$\mathfrak{c}$.\footnote{Simplicity of $\mathfrak{sp}_{4}(\RR)$ is the one
fact from Lie theory used here, and it is elementary: the complexification
$\mathfrak{sp}_{4}(\CC)$ has the irreducible root system $C_{2}$, a complex
semisimple Lie algebra with irreducible root system is simple, and an ideal of a
real form is an ideal of the complexification after tensoring with $\CC$, so
a real form of a simple complex Lie algebra is simple. Nothing finer than this
is needed, in particular no statement about the abstract group
$\mathrm{PSp}_{4}(\RR)$; and the conclusion
$\mathfrak{p}+[\mathfrak{p},\mathfrak{p}]=\mathfrak{c}$ is also checked by
direct computation in item (XXIV), at the two members of the pencil that carry
a rational point.} Hence $z$ is annihilated by $\mathfrak{c}$ and, since
$\Sp(V_{+})$ is connected, invariant under all of $\mathcal{C}_{k}$.
So the classes of type $(1,1)$ at every point of $S_{k}$ are the invariants of
$\Sp(V_{+})$ in $\bigwedge^{2}(V_{+}\otimes\RR^{2})
=\bigl(\bigwedge^{2}V_{+}\otimes S^{2}\RR^{2}\bigr)
\oplus\bigl(S^{2}V_{+}\otimes\bigwedge^{2}\RR^{2}\bigr)$, a space of dimension
$3$ spanned by $\eta$, $x_{k}$ and $x_{k}(i\cdot,\cdot)$. Item (XXIV)
confirms both steps over $\QQ$: for $0\le k\le4$ the centraliser of $i$ and
$\phi$ in $\mathfrak{sp}(V,E)$ has dimension $10$ and exactly three invariants
in $\bigwedge^{2}V^{*}$; and at the values of $k$ where $\beta(k)$ is a
rational square, $k=-1$ for $d=1$ and $k=-2$ for $d=3$, a rational point
$J_{0}$ of $S_{k}$ is written down, $\mathfrak{p}=[\mathfrak{c},J_{0}]$ has
dimension $6$, and $\mathfrak{p}+[\mathfrak{p},\mathfrak{p}]$ has dimension
$10$. Every other rational class $z$ is of type $(1,1)$ only on
a proper closed analytic subset of the connected manifold $S_{k}$, and a very
general point avoids these countably many subsets.

For the bound in (iv), write $e=\int\eta^{4}$, $\iota x=x(i\cdot,\cdot)$ and
$M_{k}=(\QQ x_{k}+\QQ\iota x_{k})\cap H^{2}(\Lambda,\ZZ)$. Since
$\phi_{\iota x}=\phi_{x}\,i$ and $i\phi_{x}=-\phi_{x}i$, the identity of (i)
gives $I((a+b\iota)x_{k})=(a^{2}+db^{2})I(x_{k})$ and
$I(x_{k})=\beta(k)e/3$. The index of $\ZZ x_{k}+\ZZ\iota x_{k}$ in $M_{k}$
divides the gcd $G(k)$ of the $2\times2$ minors of the matrix with rows
$x_{k}$, $\iota x_{k}$. These minors are integer polynomials of degree at most
two in $k$, so $G(k)$ divides every element of the ideal they generate in
$\ZZ[k]$ that is a constant; item (XXIV) finds such constants, resultants of
pairs of minors, whose gcd is $N_{d}$. Hence $N_{d}M_{k}\subset\ZZ x_{k}+\ZZ\iota
x_{k}$, and every nonzero $w\in M_{k}$ has $I(w)\ge I(x_{k})/N_{d}^{2}$. Now let
$y\in\NS\setminus\QQ\eta$ and $z=y-t(y)\eta$. Then $0\ne ez\in M_{k}$, because
$e\,y$ and $e\,t(y)=\int y\eta^{3}$ are integral, and
$I(y)=t(y)^{2}e+I(z)\ge I(ez)/e^{2}\ge\beta(k)/(3N_{d}^{2}e)$. The final
assertion follows from \Cref{lem:nefnorm}(ii).
\end{proof}

The bound in (iv) is far from sharp. The exact values computed in item (XXIV)
for $0\le k\le10$ are $\mu_{k}=\beta(k)\int\eta^{4}/3$ or half of it for
$d=1$, and $\mu_{k}=\beta(k)\int\eta^{4}/12$ for $d=3$; the proof loses the
factor $(\int\eta^{4})^{2}$ in passing from $y$ to its primitive part, and
only the growth in $k$ matters here. Each
locus of the pencil carries algebraic Weil classes, by the divisorial
criterion, and on the $k$-th of them the divisors that produce those classes
have degree growing like $k$. This is \Cref{prop:divroute} seen on explicit
quaternionic loci in one fixed polarised lattice; \Cref{fig:pencil} draws
it.

\begin{figure}[t]
\centering
  \includegraphics[width=\textwidth]{fig_pencil}
\caption{\Cref{prop:pencil} at $d=3$. Left: the members $S_{k}$ of the
pencil, each a Siegel threefold inside the four-dimensional family, with the
member $k=-2$, where $\beta=1$ and a rational point $J_{0}$ can be written
down, drawn apart. Right: along the pencil the ratio $\mu_{k}/\beta(k)$ is
the constant $2592$, so the Weil class is the same size at every member, while
the least degree of a divisor not proportional to $\eta$ grows linearly in
$k$. The class does not change; the divisors that produce it must.}
\label{fig:pencil}
\end{figure}

\begin{remark}[At $n=2$ the hypothesis is weaker still]\label{rem:ntwo}
Let $n=2$, so $\dim\cS=n^{2}=4$. The quaternionic locus of
\Cref{setup:quat} is the symmetric space of $\Sp(4,\RR)$ inside
$\cD_{2,2}$,\footnote{Left multiplication by the quaternion algebra
$B=(-d,b)$ commutes with $K$, and $B\otimes\RR\cong M_{2}(\RR)$ because
$b>0$, so Morita equivalence writes $V_{\RR}$ as $W\otimes_{\RR}\RR^{2}$ with
$W$ symplectic of dimension $2n$. The complex structures commuting with $B$
are then the complex structures on $W$, that is the points of the Siegel
space $\mathfrak{H}_{n}$, of dimension $n(n+1)/2$. The codimension in
$\cD_{n,n}$ is $n^{2}-n(n+1)/2=n(n-1)/2$, which is the count verified in
item (XII) of \Cref{sec:verification} for $n\le10$ and in
\texttt{quat\_locus\_codimension} of \Cref{sec:lean}.} of dimension
$n(n+1)/2=3$, so it is a divisor. Its $G(\QQ)$-translates are infinitely many
pairwise distinct irreducible closed subvarieties of $\cS$ of dimension $3$,
since a finite union of proper closed analytic subsets cannot be dense. If
infinitely many of them carry cycles with one and the same data
$(\underline{P},\underline{q})$, then the closed set
$\Sigma_{\underline{P},\underline{q}}$ has an irreducible component containing
infinitely many distinct irreducible threefolds, so that component has
dimension at least $4$ and therefore equals $\cS$. No density is needed at
$n=2$, only infinitude. Combined with the theorem of Moonen and Zarhin
\cite{MZ99} that every Hodge class on an abelian fourfold lies in the
subalgebra generated by divisor classes and Weil classes, the conclusion would
be the Hodge conjecture for all abelian fourfolds of Weil type, hence for all
abelian fourfolds. That conclusion is now Markman's theorem
\cite[Corollary 1.6.1]{Mar25a}, reached by a different route; at $n=2$ the
argument above would be a second proof, and it is at $n\ge3$ that it has
something to add.
\end{remark}
